% Part: normal-modal-logic
% Chapter: axioms-systems
% Section: provability-properties

\documentclass[../../../include/open-logic-section]{subfiles}

\begin{document}

\olfileid[pt]{nml}{prf}{prp}

\olsection{Propriedades da \usetoken{S}{derivability}}

\begin{prop}\ollabel{prop:derivabilityfacts}
  Seja $\Sigma$ um sistema modal e $\Gamma$ um conjunto de
  !!{formula}s modais. Valem as seguintes propriedades:
  \begin{enumerate}
  \item\ollabel{prop:derivabilityfacts-monotonicity} \emph{Monotonicidade}: Se
    $\Gamma \Proves[\Sigma] !A$ e $\Gamma \subseteq \Delta$, então
    $\Delta \Proves[\Sigma] !A$;
  \item\ollabel{prop:derivabilityfacts-reflexivity}
    \emph{Reflexividade}: Se $!A \in \Gamma$, então $\Gamma
    \Proves[\Sigma] !A$;
  \item\ollabel{prop:derivabilityfacts-cut} \emph{Corte}: Se $\Gamma
    \Proves[\Sigma] !A$ e $\Delta \cup \{!A\} \Proves[\Sigma] !B$,
    então $\Gamma \cup \Delta \Proves[\Sigma] !B$;
  \item\ollabel{prop:derivabilityfacts-deduction} \emph{Teorema da dedução}: $\Gamma \cup \{!B\}
    \Proves[\Sigma] !A$ se e somente se $\Gamma \Proves[\Sigma] !B \lif
      !A$;
  \item \ollabel{prop:derivabilityfacts-ruleT}% \emph{Regra T}: Se
    $\Gamma \Proves[\Sigma] !A_1$ e \dots e $\Gamma
    \Proves[\Sigma] !A_n$ e $!A_1 \to (!A_2 \lif \cdots (!A_n \lif
    !B)\cdots)$ é uma instância tautológica, então $\Gamma
    \Proves[\Sigma] !B$.
  \end{enumerate}
\end{prop}

A demonstração é um exercício fácil. A parte
\olref{prop:derivabilityfacts-ruleT} de \olref{prop:derivabilityfacts}
nos dá, por exemplo, que, se $\Gamma \Proves[\Sigma] !A \lor !B$ e
$\Gamma \Proves[\Sigma] \lnot !A$, então $\Gamma \Proves[\Sigma]
!B$. Além disso, daqui em diante escrevemos
$\Gamma, !A \Proves[\Sigma] !B$ em vez de
$\Gamma \cup \{ !A\} \Proves[\Sigma] !B$.

\begin{defn}
  Um conjunto $\Gamma$ é \emph{dedutivamente fechado} relativamente a
  um sistema $\Sigma$ se e somente se $\Gamma \Proves[\Sigma] !A$ implica
  $!A \in \Gamma$.
\end{defn}

\end{document}
